% Options for packages loaded elsewhere
\PassOptionsToPackage{unicode}{hyperref}
\PassOptionsToPackage{hyphens}{url}
\documentclass[
  ignorenonframetext,
]{beamer}
\newif\ifbibliography
\usepackage{pgfpages}
\setbeamertemplate{caption}[numbered]
\setbeamertemplate{caption label separator}{: }
\setbeamercolor{caption name}{fg=normal text.fg}
\beamertemplatenavigationsymbolsempty
% remove section numbering
\setbeamertemplate{part page}{
  \centering
  \begin{beamercolorbox}[sep=16pt,center]{part title}
    \usebeamerfont{part title}\insertpart\par
  \end{beamercolorbox}
}
\setbeamertemplate{section page}{
  \centering
  \begin{beamercolorbox}[sep=12pt,center]{section title}
    \usebeamerfont{section title}\insertsection\par
  \end{beamercolorbox}
}
\setbeamertemplate{subsection page}{
  \centering
  \begin{beamercolorbox}[sep=8pt,center]{subsection title}
    \usebeamerfont{subsection title}\insertsubsection\par
  \end{beamercolorbox}
}
% Prevent slide breaks in the middle of a paragraph
\widowpenalties 1 10000
\raggedbottom
\AtBeginPart{
  \frame{\partpage}
}
\AtBeginSection{
  \ifbibliography
  \else
    \frame{\sectionpage}
  \fi
}
\AtBeginSubsection{
  \frame{\subsectionpage}
}
\usepackage{iftex}
\ifPDFTeX
  \usepackage[T1]{fontenc}
  \usepackage[utf8]{inputenc}
  \usepackage{textcomp} % provide euro and other symbols
\else % if luatex or xetex
  \usepackage{unicode-math} % this also loads fontspec
  \defaultfontfeatures{Scale=MatchLowercase}
  \defaultfontfeatures[\rmfamily]{Ligatures=TeX,Scale=1}
\fi
\usepackage{lmodern}
\ifPDFTeX\else
  % xetex/luatex font selection
\fi
% Use upquote if available, for straight quotes in verbatim environments
\IfFileExists{upquote.sty}{\usepackage{upquote}}{}
\IfFileExists{microtype.sty}{% use microtype if available
  \usepackage[]{microtype}
  \UseMicrotypeSet[protrusion]{basicmath} % disable protrusion for tt fonts
}{}
\makeatletter
\@ifundefined{KOMAClassName}{% if non-KOMA class
  \IfFileExists{parskip.sty}{%
    \usepackage{parskip}
  }{% else
    \setlength{\parindent}{0pt}
    \setlength{\parskip}{6pt plus 2pt minus 1pt}}
}{% if KOMA class
  \KOMAoptions{parskip=half}}
\makeatother
\usepackage{color}
\usepackage{fancyvrb}
\newcommand{\VerbBar}{|}
\newcommand{\VERB}{\Verb[commandchars=\\\{\}]}
\DefineVerbatimEnvironment{Highlighting}{Verbatim}{commandchars=\\\{\}}
% Add ',fontsize=\small' for more characters per line
\newenvironment{Shaded}{}{}
\newcommand{\AlertTok}[1]{\textcolor[rgb]{1.00,0.00,0.00}{\textbf{#1}}}
\newcommand{\AnnotationTok}[1]{\textcolor[rgb]{0.38,0.63,0.69}{\textbf{\textit{#1}}}}
\newcommand{\AttributeTok}[1]{\textcolor[rgb]{0.49,0.56,0.16}{#1}}
\newcommand{\BaseNTok}[1]{\textcolor[rgb]{0.25,0.63,0.44}{#1}}
\newcommand{\BuiltInTok}[1]{\textcolor[rgb]{0.00,0.50,0.00}{#1}}
\newcommand{\CharTok}[1]{\textcolor[rgb]{0.25,0.44,0.63}{#1}}
\newcommand{\CommentTok}[1]{\textcolor[rgb]{0.38,0.63,0.69}{\textit{#1}}}
\newcommand{\CommentVarTok}[1]{\textcolor[rgb]{0.38,0.63,0.69}{\textbf{\textit{#1}}}}
\newcommand{\ConstantTok}[1]{\textcolor[rgb]{0.53,0.00,0.00}{#1}}
\newcommand{\ControlFlowTok}[1]{\textcolor[rgb]{0.00,0.44,0.13}{\textbf{#1}}}
\newcommand{\DataTypeTok}[1]{\textcolor[rgb]{0.56,0.13,0.00}{#1}}
\newcommand{\DecValTok}[1]{\textcolor[rgb]{0.25,0.63,0.44}{#1}}
\newcommand{\DocumentationTok}[1]{\textcolor[rgb]{0.73,0.13,0.13}{\textit{#1}}}
\newcommand{\ErrorTok}[1]{\textcolor[rgb]{1.00,0.00,0.00}{\textbf{#1}}}
\newcommand{\ExtensionTok}[1]{#1}
\newcommand{\FloatTok}[1]{\textcolor[rgb]{0.25,0.63,0.44}{#1}}
\newcommand{\FunctionTok}[1]{\textcolor[rgb]{0.02,0.16,0.49}{#1}}
\newcommand{\ImportTok}[1]{\textcolor[rgb]{0.00,0.50,0.00}{\textbf{#1}}}
\newcommand{\InformationTok}[1]{\textcolor[rgb]{0.38,0.63,0.69}{\textbf{\textit{#1}}}}
\newcommand{\KeywordTok}[1]{\textcolor[rgb]{0.00,0.44,0.13}{\textbf{#1}}}
\newcommand{\NormalTok}[1]{#1}
\newcommand{\OperatorTok}[1]{\textcolor[rgb]{0.40,0.40,0.40}{#1}}
\newcommand{\OtherTok}[1]{\textcolor[rgb]{0.00,0.44,0.13}{#1}}
\newcommand{\PreprocessorTok}[1]{\textcolor[rgb]{0.74,0.48,0.00}{#1}}
\newcommand{\RegionMarkerTok}[1]{#1}
\newcommand{\SpecialCharTok}[1]{\textcolor[rgb]{0.25,0.44,0.63}{#1}}
\newcommand{\SpecialStringTok}[1]{\textcolor[rgb]{0.73,0.40,0.53}{#1}}
\newcommand{\StringTok}[1]{\textcolor[rgb]{0.25,0.44,0.63}{#1}}
\newcommand{\VariableTok}[1]{\textcolor[rgb]{0.10,0.09,0.49}{#1}}
\newcommand{\VerbatimStringTok}[1]{\textcolor[rgb]{0.25,0.44,0.63}{#1}}
\newcommand{\WarningTok}[1]{\textcolor[rgb]{0.38,0.63,0.69}{\textbf{\textit{#1}}}}
\setlength{\emergencystretch}{3em} % prevent overfull lines
\providecommand{\tightlist}{%
  \setlength{\itemsep}{0pt}\setlength{\parskip}{0pt}}
% Slide layout defaults for warning-clean scientific decks.
\usepackage{etoolbox}
\IfFileExists{xurl.sty}{\usepackage{xurl}}{}
\IfFileExists{seqsplit.sty}{\usepackage{seqsplit}}{\newcommand{\seqsplit}[1]{#1}}
\protected\def\breakseq#1{\seqsplit{#1}}
\protected\def\breaktt#1{\begingroup\ttfamily\seqsplit{#1}\endgroup}
\setlength{\emergencystretch}{6em}
\tolerance=5000
\hbadness=10000
\hfuzz=1pt
\setlength{\tabcolsep}{2pt}
\AtBeginEnvironment{longtable}{\tiny\renewcommand{\arraystretch}{0.86}\setlength{\tabcolsep}{1pt}}
\AtBeginEnvironment{tabular}{\tiny\renewcommand{\arraystretch}{0.86}\setlength{\tabcolsep}{1pt}}

% Natbib and cross-reference fallbacks — slides don't load natbib
% or cleveref, but combined-PDF manuscript prose may emit these
% commands. The fallback renders citations as a bracketed key list
% and cross-references as detokenized labels so slides don't fail on
% undefined control sequences or raw underscores. \providecommand is
% a no-op if packages load later.
\providecommand{\citep}[1]{[#1]}
\providecommand{\citet}[1]{#1}
\providecommand{\citealp}[1]{#1}
\providecommand{\citeauthor}[1]{#1}
\providecommand{\citeyear}[1]{#1}
\providecommand{\cref}[1]{\texttt{\detokenize{#1}}}
\providecommand{\Cref}[1]{\texttt{\detokenize{#1}}}
\usepackage{bookmark}
\IfFileExists{xurl.sty}{\usepackage{xurl}}{} % add URL line breaks if available
\urlstyle{same}
% fallback for those not using the hyperref driver hyperxmp:
\makeatletter
\@ifundefined{xmpquote}{\newcommand{\xmpquote}[1]{#1}}{}
\makeatother
\hypersetup{
  hidelinks,
  pdfcreator={LaTeX via pandoc}}

\author{\texorpdfstring{}{}}
\date{}

\begin{document}

\begin{frame}[fragile,allowframebreaks]{Reproducibility}
\protect\phantomsection\label{reproducibility}
Reproducibility is not a claim this manuscript makes about itself; it is
a property the code enforces on every run. Every number that appears
below the Determinism heading --- f84a8f9dbcb18e37, 42, 493, 96.28 ---
is a token substituted at render time by
\breaktt{src/manuscript\_variables.py::generate\_variables()} from a live
grammar load and a live pytest run
(\breaktt{src/manuscript\_variables.py::measure\_test\_summary()}).
Neither function accepts a hardcoded literal as a fallback: if the
subprocess pytest run cannot be parsed, \breaktt{measure\_test\_summary}
returns the string \texttt{"pending"} for both the test count and the
coverage percentage rather than a plausible-looking number
\breakseq{[@reproducible\_builds]}. This design choice is a direct consequence
of an earlier failure mode in this project --- a manuscript draft that
stated fixed values (``Tests: 371 · Coverage: 99.94\%'') in prose
instead of through the token pipeline. A number that cannot be traced
back to a specific function call in \texttt{src/} is, for the purposes
of this project, not a number this manuscript is permitted to assert.

\begin{block}{Determinism guarantee}
\protect\phantomsection\label{determinism-guarantee}
Given the same \breaktt{manuscript/config.yaml} grammar definition
(grammar hash \breaktt{f84a8f9dbcb18e37}), the same seed (\texttt{42}),
and the same \breaktt{template\_autopoiesis} source tree,
\breaktt{scripts/autopoiesis.py\ expand} produces byte-identical
selections on every invocation, and \texttt{materialize} produces a
byte-identical child project tree. The determinism chain has three
concrete steps, each implemented as a pure function with no random or
wall-clock input:

\begin{enumerate}
\tightlist
\item
  \textbf{Selection.} For every grammar slot,
  \breaktt{src/expand.py::\_digest\_index()} builds the key
  \texttt{f"\{seed\}\textbackslash{}x1f\{slot\_name\}\textbackslash{}x1f\{ordinal\}\textbackslash{}x1f\{\textquotesingle{},\textquotesingle{}.join(options)\}"}
  (a unit separator, \texttt{\textbackslash{}x1f}, joins the fields so
  that no combination of seed/name/ ordinal/option values can collide
  across a field boundary), hashes it with SHA-256, and reduces the
  first eight bytes of the digest modulo \texttt{len(options)} to obtain
  the chosen option's index. Because the digest is a pure function of
  \breaktt{(seed,\ slot\_name,\ ordinal,\ options)}, the same grammar and
  seed walk every slot to the same choice on every run --- there is no
  \texttt{random.choice}, no \texttt{numpy.random}, and no
  seed-independent entropy source anywhere in the selection path.
\item
  \textbf{Spec identity.} The resolved selections, together with the
  grammar hash and seed, are assembled into a \texttt{Spec} dataclass
  (\breaktt{src/expand.py::Spec}). Its \texttt{spec\_hash} property
  serializes the full spec to canonical JSON (\texttt{sort\_keys=True},
  compact separators) and takes the first sixteen hex characters of the
  SHA-256 of that string. Sorting keys before hashing means insertion
  order in the underlying dict cannot perturb the hash --- only the
  actual selections can.
\item
  \textbf{Tree identity.} \texttt{materialize()}
  (\breaktt{src/materialize.py}) writes every vendored and generated file
  into the child project directory, then folds the complete
  \texttt{\{relative\_path:\ content\}} mapping through
  \breaktt{src/integrity.py::tree\_hash\_from\_content\_hashes()}. That
  function sorts the mapping lexicographically by path before hashing
  (\texttt{"\textbackslash{}n".join(f"\{k\}:\{v\}"\ for\ k,\ v\ in\ sorted(...))})
  specifically so that filesystem iteration order --- which is not
  stable across platforms or \texttt{dict} construction paths --- cannot
  change the result. The hash is computed from file \emph{contents}, not
  from file metadata (mtime, permissions, inode), so copying a child
  project to a new machine or re-materializing it a year later
  reproduces the same tree hash as long as the source and the seed are
  unchanged.
\end{enumerate}

Multiple children can be derived from one root seed without collisions:
given a \texttt{base\_seed} and an integer \texttt{index},
\breaktt{src/expand.py::derive\_seed()} hashes
\texttt{f"\{base\_seed\}\textbackslash{}x1f\{index\}"} through SHA-256
and folds the first eight bytes to a new integer seed.
\breaktt{sample(grammar,\ count)} calls \texttt{expand()} once per
derived seed, so a batch of \texttt{count} children is itself
deterministic --- the same \texttt{base\_seed} and \texttt{count} yield
the same sequence of child seeds on every invocation, hence the same
sequence of children.
\end{block}

\begin{block}{The seal.json provenance mechanism}
\protect\phantomsection\label{the-seal.json-provenance-mechanism}
Materialization and sealing are deliberately two separate steps, and the
second one is optional. \texttt{materialize()} writes a
\texttt{provenance.json} into the child root containing the schema
version
(\breaktt{PROVENANCE\_SCHEMA\_VERSION\ =\ "autopoiesis/provenance/1"}),
the full resolved spec (\texttt{spec.to\_dict()}), the tree hash, and
the sorted list of every file path that was written. This is the record
\texttt{verify\_child()} (\texttt{src/verify.py}) recomputes against
later.

A child project can additionally be \textbf{sealed}:
\breaktt{scripts/seal\_child.py} (and the pipeline-facing wrapper
\breaktt{scripts/04\_seal.py}, which seals the most-recently materialized
child under \breaktt{output/children/}) reads \texttt{provenance.json},
re-runs \texttt{verify\_child()} against the live files as a pre-seal
sanity check, and writes \texttt{seal.json} alongside it. The seal
payload --- built by \breaktt{src/sealing.py::build\_payload()} --- is a
compact JSON object
\texttt{\{"spec\_hash":\ ...,\ "tree\_hash":\ ...,\ "seed":\ ...\}}. A
shorter, colon-delimited variant (\breaktt{build\_barcode\_payload()},
truncating each hash to its first eight hex characters) exists for
embedding into a QR code or barcode image via
\breaktt{src/sealing.py::qr\_matrix()} / \texttt{qr\_image()}, so that a
printed or exported artifact can carry a scannable, self-describing
pointer back to the exact spec and tree hash that produced it. Both the
QR encoder and its optional decode path (\breaktt{read\_qr\_matrix()},
which depends on \texttt{pyzbar} and \texttt{PIL}) degrade gracefully
when those optional dependencies are absent: \texttt{qr\_matrix()} falls
back to a deterministic checkerboard stub so callers and tests do not
hard-fail on a missing image dependency, and \breaktt{read\_qr\_matrix()}
simply returns an empty string. The seal itself does not gate
materialization --- a child project is fully valid and independently
verifiable from \texttt{provenance.json} alone; \texttt{seal.json} is an
additive, portable pointer, not a second source of truth. Sealing does
not currently run inside \breaktt{verify\_child\_full()}; it is a
separate check (\breaktt{src/verify.py::verify\_seal()}) invoked only
when a caller explicitly asks whether a seal exists, parses, and carries
a \texttt{spec\_hash}.
\end{block}

\begin{block}{SHA-256 vs.~Merkle integrity profiles}
\protect\phantomsection\label{sha-256-vs.-merkle-integrity-profiles}
\breaktt{src/integrity.py} provides two distinct ways of turning a
collection of hashes into one summary digest, and the project does not
conflate them:

\begin{itemize}
\tightlist
\item
  \textbf{\breaktt{tree\_hash\_from\_content\_hashes()}} --- the profile
  used by \texttt{materialize()}/\texttt{verify\_child()} above --- is a
  single-pass, order-independent digest over a
  \texttt{\{path:\ content\}} mapping. It answers exactly one question:
  does this set of files, at these paths, have these exact contents? It
  is cheap to compute and cheap to verify, but a single changed file
  forces a full recompute over every path to detect \emph{which} file
  changed --- the digest itself carries no positional structure.
\item
  \textbf{\texttt{merkle\_root()}} builds a binary Merkle tree
  \breakseq{[@merkle\_tree\_provenance]} over an \emph{ordered} list of hex
  digests: each level pairwise-concatenates adjacent nodes and hashes
  the concatenation, promoting an unpaired trailing node unchanged to
  the next level, until one root remains (the empty list is defined as
  the SHA-256 of the empty string, not a magic sentinel). Because the
  tree is addressable node-by-node, a Merkle profile supports proving
  that one specific file's hash is part of the committed set without
  re-hashing every other file --- a property the flat tree-hash profile
  does not have.
\end{itemize}

Both are exercised in this codebase, but at present the
materialize/verify path in \breaktt{src/materialize.py} and
\texttt{src/verify.py} uses the flat, order-independent tree hash
exclusively; \texttt{merkle\_root()} is available in
\breaktt{src/integrity.py} and covered by its own tests as an independent
integrity primitive, not yet as the provenance root written into
\texttt{provenance.json}. A manuscript describing this project should
not claim Merkle-tree provenance for \texttt{provenance.json} today ---
that would be exactly the kind of prose-outruns-code gap this project's
honesty checks exist to catch (\texttt{src/honesty.py}).
\end{block}

\begin{block}{Recompute / verify workflow}
\protect\phantomsection\label{recompute-verify-workflow}
\begin{Shaded}
\begin{Highlighting}[]
\CommentTok{\# 1. Expand the grammar into a resolved spec (pure function of seed + config)}
\ExtensionTok{uv}\NormalTok{ run python scripts/autopoiesis.py expand }\AttributeTok{{-}{-}seed}\NormalTok{ 42 }\AttributeTok{{-}{-}output}\NormalTok{ output/spec.json}

\CommentTok{\# 2. Materialize the spec into a runnable child project + provenance.json}
\ExtensionTok{uv}\NormalTok{ run python scripts/autopoiesis.py materialize }\AttributeTok{{-}{-}seed}\NormalTok{ 42 }\AttributeTok{{-}{-}out{-}root}\NormalTok{ output/children}

\CommentTok{\# 3. Recompute the tree hash from the live files and compare to provenance.json}
\ExtensionTok{uv}\NormalTok{ run python scripts/autopoiesis.py verify output/children/}\OperatorTok{\textless{}}\NormalTok{child\_name}\OperatorTok{\textgreater{}}

\CommentTok{\# 4. (optional) Seal the child — writes seal.json next to provenance.json}
\ExtensionTok{uv}\NormalTok{ run python scripts/seal\_child.py output/children/}\OperatorTok{\textless{}}\NormalTok{child\_name}\OperatorTok{\textgreater{}}
\end{Highlighting}
\end{Shaded}

\texttt{verify} (\breaktt{src/cli.py::cmd\_verify} →
\breaktt{src/verify.py::verify\_child\_full()}) performs four checks in
sequence and exits non-zero if any fails: \breaktt{provenance\_exists},
\breaktt{provenance\_parseable}, \breaktt{all\_files\_present} (every path
listed in \texttt{provenance.json{[}"files"{]}} still exists on disk),
and \breaktt{tree\_hash\_matches} (the hash recomputed from the live file
contents equals the hash recorded at materialization time), followed by
a \breaktt{schema\_version\_correct} check against the constant
\breaktt{PROVENANCE\_SCHEMA\_VERSION}. There is no partial-credit
outcome: a single missing file or a single byte of drift in any tracked
file flips \breaktt{tree\_hash\_matches} to \texttt{False}, because the
tree hash is a single SHA-256 over the sorted, concatenated
\texttt{path:content} pairs --- there is no per-file tolerance to fall
back on.

Re-running steps 1--2 with the same seed and the same source tree on the
same machine reproduces the identical spec hash and tree hash reported
by step 3; this is the operational meaning of ``deterministic'' for this
project, and it is exactly what a reviewer or a downstream user can
check for themselves without trusting any claim made in this document.
\end{block}

\begin{block}{Toolchain}
\protect\phantomsection\label{toolchain}
\begin{itemize}
\tightlist
\item
  Python ≥ 3.10
\item
  \texttt{numpy}, \texttt{matplotlib}, \texttt{pyyaml} (runtime)
\item
  \texttt{pytest}, \texttt{pytest-cov} (test)
\item
  \texttt{hypothesis} \breakseq{[@claessen2000quickcheck;
  @maciver2019hypothesis]} (declared under this project's \texttt{dev}
  optional-dependency group in \texttt{pyproject.toml}, but imported
  unconditionally at the top of \breaktt{test\_property\_invariants.py}
  --- a real dependency of that test module, not a soft,
  try/except-guarded one; see Methods)
\item
  \texttt{qrcode}, \texttt{pillow} (optional ---
  \texttt{src/sealing.py}'s QR/barcode payload encoding degrades to a
  deterministic stub when unavailable; \texttt{pyzbar} is needed only
  for the optional decode path)
\end{itemize}

Outside the parent \texttt{template} repository checkout,
\texttt{STANDALONE.md} documents what is and is not vendored into a
generated child project: single-file modules under \texttt{primitives/}
and \texttt{figures.py} are copied in and the child runs entirely from
its own vendored sources at runtime, with no dependency on the parent's
\texttt{src/}; optional infrastructure modules (logging, glossary
generation, figure management) are vendored as single-file stubs when
the corresponding module cannot be found under the parent repo's
\texttt{infrastructure/} tree. PDF/HTML rendering and prerender
validation are explicitly \textbf{not} vendored --- a standalone child
can be expanded, materialized, tested, and verified, but manuscript
rendering to PDF still requires the parent repository's
\breaktt{infrastructure/rendering} and \breaktt{infrastructure/validation}
modules. Byte-stability of a materialized tree is documented as
within-platform only (same Python interpreter, same OS) rather than an
unconditional cross-platform guarantee.
\end{block}

\begin{block}{Build command}
\protect\phantomsection\label{build-command}
\begin{Shaded}
\begin{Highlighting}[]
\ExtensionTok{uv}\NormalTok{ run pytest projects/templates/template\_autopoiesis/tests/ }\DataTypeTok{\textbackslash{}}
  \AttributeTok{{-}{-}cov}\OperatorTok{=}\NormalTok{projects/templates/template\_autopoiesis/src }\DataTypeTok{\textbackslash{}}
  \AttributeTok{{-}{-}cov{-}fail{-}under}\OperatorTok{=}\NormalTok{90}
\end{Highlighting}
\end{Shaded}

This is the same command family \breaktt{measure\_test\_summary()} runs
internally (with \texttt{-\/-cov-branch} added explicitly, matching the
repo-root \texttt{pyproject.toml}'s \texttt{branch\ =\ true} setting
rather than this project's own, so the reported 96.28 agrees with the
authoritative CI gate methodology rather than silently using a different
one) to produce the 493 and 96.28 tokens substituted above --- the same
build, not a paraphrase of it, is the source of both this document's
prose and its own regeneration.
\end{block}
\end{frame}

\end{document}
